%%% Packages included in thesis-umich class
%\RequirePackage[margin=1in,footskip=8pt,headsep=0.4cm,headheight=\baselineskip]{geometry}
%\RequirePackage{amsmath}
%\RequirePackage{amsfonts}
%\RequirePackage{amssymb}
%\RequirePackage{graphicx}
%\RequirePackage{subcaption}
%\RequirePackage{times}
%\RequirePackage{natbib}
%\RequirePackage{verbatim}
%\RequirePackage{upquote}
%\RequirePackage{textcomp}
%\RequirePackage{setspace}
%\RequirePackage{ifthen}
%\RequirePackage{soul}
%\RequirePackage{float}
%\RequirePackage{acronym}
%\RequirePackage{makeidx}
%\RequirePackage{fancyhdr}
%\RequirePackage{multicol}

% Include your packages here by including a \usepackage{<package_name>} command.
\usepackage{blindtext} % Example package which populates the ever-popular lorem ipsum text.
\usepackage{listings}
% ---- Code listings (added 2026-11-12) -------------------------------------------------------
% Style from the Overleaf code-listing example.
\usepackage{xcolor}
\definecolor{codegreen}{rgb}{0,0.6,0}
\definecolor{codegray}{rgb}{0.5,0.5,0.5}
\definecolor{codepurple}{rgb}{0.58,0,0.82}
\definecolor{backcolour}{rgb}{0.95,0.95,0.92}
\lstdefinestyle{mystyle}{
  backgroundcolor=\color{backcolour},   commentstyle=\color{codegreen},
  keywordstyle=\color{magenta},
  numberstyle=\tiny\color{codegray},
  stringstyle=\color{codepurple},
  basicstyle=\ttfamily\footnotesize,
  breakatwhitespace=false,
  breaklines=true,
  captionpos=b,
  keepspaces=true,
  numbers=left,
  numbersep=5pt,
  showspaces=false,
  showstringspaces=false,
  showtabs=false,
  tabsize=2
}
\lstset{style=mystyle}

% \codefile[language]{path}{caption}{label}
% A code file printed in full. Unlike the class's "program" float it can run across pages,
% and it still gets a number and a line in the List of Programs.
\newcommand{\codefile}[4][]{%
  \refstepcounter{program}\label{#4}%
  \addcontentsline{lop}{program}{\protect\numberline{\theprogram}#3}%
  \lstinputlisting[language={#1},title={Program \theprogram: #3}]{#2}}

% --- Vega Fluo: citations link to the chapter's own reference list (generated in .margin/chapter-refs) ---
\RequirePackage{currfile}
\makeatletter
\InputIfFileExists{.margin/chapter-refs/_index.tex}{}{}
\AtBeginDocument{\@ifundefined{hyper@natlinkstart}{}{%
  \let\vf@natlinkstart\hyper@natlinkstart
  \let\vf@natlinkbreak\hyper@natlinkbreak
  \def\hyper@natlinkstart#1{\ifcsname vf@list@\currfilebase\endcsname\Hy@backout{#1}\hyper@linkstart{cite}{cite.\currfilebase.#1}\def\hyper@nat@current{#1}\else\vf@natlinkstart{#1}\fi}%
  \def\hyper@natlinkbreak#1#2{\ifcsname vf@list@\currfilebase\endcsname\hyper@linkend#1\hyper@linkstart{cite}{cite.\currfilebase.#2}\else\vf@natlinkbreak{#1}{#2}\fi}}}
\setcitestyle{authoryear,round}
\makeatother
% --- end Vega Fluo ---

% --- Vega Fluo: a placeholder box for a figure that is not final yet: \figureplaceholder[height]{what goes here} ---
\providecommand{\figureplaceholder}[2][2.4in]{\fbox{\parbox[c][#1][c]{0.85\linewidth}{\centering\sffamily\small Figure placeholder\\[0.6ex]\footnotesize #2}}}
% --- end Vega Fluo ---

% --- Vega Fluo: figure versions. \figureversion{text} uses each figure file as named; \figureversion{clean}
% uses figures/<name>_clean.pdf instead wherever that file exists (statistics then go in the legend and Results).
\makeatletter
\newif\ifvf@cleanfigs
\newcommand{\figureversion}[1]{\def\vf@tmp{#1}\def\vf@clean{clean}\ifx\vf@tmp\vf@clean\vf@cleanfigstrue\else\vf@cleanfigsfalse\fi}
\def\vf@trypdf#1.pdf#2\@nil{\IfFileExists{#1_clean.pdf}{\def\vf@file{#1_clean.pdf}}{}}
\newcommand{\vf@includegraphics}[2][]{\def\vf@file{#2}\ifvf@cleanfigs\vf@trypdf#2.pdf\@nil\fi\vf@origgraphics[#1]{\vf@file}}
\AtBeginDocument{\let\vf@origgraphics\includegraphics\renewcommand{\includegraphics}{\@ifstar{\vf@origgraphics*}{\vf@includegraphics}}}
\makeatother
